\scanpage{8}
Soit $\chi:T\mathbin{-}G_m$
\[
\begin{aligned}
&T_\chi=(\operatorname{Ker}\chi)^\circ\subset T\\
&Z_\chi=\mathrm{Cent}(T_\chi)\\
&Z_\chi/T_\chi\quad\text{a comme tore maximal }T/T_\chi\\
&\hspace{7em}\text{de dim }1
\end{aligned}
\]

\diagram{\begin{tikzcd}[ampersand replacement=\&]
{} \& \unsure{\mathrm{Psaut}|\mathrm{Birat}|\mathrm{Baut}}(T/T_\chi)\\
Z_\chi/T_\chi\arrow[ur,dashed]\arrow[r,hook] \& {\mathrm{Aut}(\underline{\mathbb P}^{1})=\mathrm{GPL}(2)}\arrow[u,hook]
\end{tikzcd}}
\inserted{\unsure{car|cas} $P^1$ de dim 1}

4 possibilités pour ss-gpe de $\unsure{\mathrm{GBL}(2)|\mathrm{GPL}(2)|\mathrm{PGL}(2)}$ \unsure{contt} tore maximal
\diagram{\(\left\{\begin{array}{ll}
\text{tore maximal de }\mathrm{GPL}(2) &: Z_\chi=T\\
\text{son normalisateur} &: Z_\chi/T\ \text{de rang }2\\
\text{Borel} &\dim Z_\chi/T=1\\
\text{tout} &Z_\chi\ \text{réductif de rg}\\
&\boxed{1/2\ \text{simple }\underline{1}}
\end{array}\right.\)}

\uline{Conclusion} Supposons $\chi\in R$, les deux premiers \unsure{son|sont} exclus.

\textcircled{1} On trouve
\[
\begin{aligned}
&\left\{\begin{aligned}
 &\dim\mathfrak g_\alpha=1\\
 &\exists !\ x_\alpha:G_a\longrightarrow Z_\alpha\subset G,\qquad
 x_\alpha(G_a)=U_\alpha,\ \mathrm{Lie}(U_\alpha)=\mathfrak g_\alpha\\
 &\hspace{12em}t x_\alpha(\lambda)t^{-1}=x_\alpha(\alpha(t)\lambda)
\end{aligned}\right.\\
&\hspace{3em}\rho_\alpha:G_m\mathbin{-}T\qquad f(x_\alpha(\lambda))=\rho_\alpha(1+\lambda),\quad
 \langle\rho_\alpha,\alpha\rangle=1
\end{aligned}
\]

\[
\underline{\mathrm{Lie}}\, Z_\chi=\mathfrak g^{0}+
 \sum_{\beta\in\unsure{Q|\mathbb Q}\chi\mathbin{\wedge}R}\mathfrak g^\beta .
\]

\textcircled{2} Si $\lambda\alpha\in R$ \quad $\lambda\neq1$, alors $\lambda=-1$ et $Z_\alpha$ est réductif de rg.ss.1, dont les 2 sgpes de Borel qui $\supset T$ sont $TU_\alpha$ et $TU_{-\alpha}$

\[
\Downarrow\ \mathrm{SGA}\,3\ \underline{\mathrm{XXVI}}
\]
$\exists$ décomposition de Lévi \qquad $G^\circ=L.U$, \quad $\underset{\text{\inserted{réductif}}}{L}\supset T$

Les racines de $L$ sont les $\alpha\in R$ tels que $-\alpha\in R$.

% archive sheet number: 7
